\chapter{Rappresentazione di insiemi disgiunti mediante Foreste di insiemi Disgiunti}

\section{Struttura della foresta}
In una rappresentazione più veloce degli insiemi disgiunti, rappresentiamo gli insiemi con alberi radicati, dove:
\begin{itemize}
    \item Ogni insieme è rappresentato da un albero radicato
    \item Ogni nodo contiene un elemento
    \item Ogni nodo punta al padre e la radice a se stessa
    \item La radice contiene il rappresentante
\end{itemize}

\subsection{Esempio grafico}
\begin{figure}[htbp]
\centering
\begin{tikzpicture}[
    vertex/.style={
        circle,
        draw=black!75,
        fill=black!7,
        inner sep=0pt,
        minimum size=8mm,
        font=\large\itshape,
        line width=0.6pt
    },
    edge/.style={
        -{Stealth[length=2.8mm, width=2.1mm]},
        line width=0.85pt,
        draw=black!90,
        shorten >=1pt,
        shorten <=1pt
    },
    loop edge/.style={
        -{Stealth[length=2.8mm, width=2.1mm]},
        line width=0.85pt,
        draw=black!90,
        looseness=6,
        in=125,
        out=55
    },
    sublabel/.style={
        font=\large,
        text=black
    }
]

% ==================== PARTE (a) ====================
\begin{scope}[xshift=0cm]
    % Albero con radice c
    \node[vertex] (c) at (2.2, 4.0) {$c$};
    \node[vertex] (h) at (1.0, 2.3) {$h$};
    \node[vertex] (e) at (3.4, 2.3) {$e$};
    \node[vertex] (b) at (1.0, 0.6) {$b$};

    % Albero con radice f
    \node[vertex] (f1) at (6.2, 4.0) {$f$};
    \node[vertex] (d1) at (6.2, 2.3) {$d$};
    \node[vertex] (g1) at (6.2, 0.6) {$g$};

    % Archi albero c
    \draw[edge] (h) -- (c);
    \draw[edge] (e) -- (c);
    \draw[edge] (b) -- (h);
    \draw[loop edge] (c) to (c);

    % Archi albero f
    \draw[edge] (d1) -- (f1);
    \draw[edge] (g1) -- (d1);
    \draw[loop edge] (f1) to (f1);

    % Etichetta (a)
    \node[sublabel] at (3.9, -1.8) {(a)};
\end{scope}

% ==================== PARTE (b) ====================
\begin{scope}[xshift=10.5cm]
    % Radice f
    \node[vertex] (f2) at (3.0, 4.0) {$f$};

    % Figli di f
    \node[vertex] (c2) at (1.4, 2.3) {$c$};
    \node[vertex] (d2) at (4.6, 2.3) {$d$};

    % Figli di c e d
    \node[vertex] (h2) at (0.2, 0.6) {$h$};
    \node[vertex] (e2) at (2.6, 0.6) {$e$};
    \node[vertex] (g2) at (4.6, 0.6) {$g$};

    % Figlio di h
    \node[vertex] (b2) at (0.2, -1.1) {$b$};

    % Archi
    \draw[edge] (c2) -- (f2);
    \draw[edge] (d2) -- (f2);
    \draw[edge] (h2) -- (c2);
    \draw[edge] (e2) -- (c2);
    \draw[edge] (g2) -- (d2);
    \draw[edge] (b2) -- (h2);
    \draw[loop edge] (f2) to (f2);

    % Etichetta (b)
    \node[sublabel] at (3.0, -1.8) {(b)};
\end{scope}

\end{tikzpicture}
\end{figure}

\section{Pseudocodice Metodi Stock}


\subsection{Make-Set}
\begin{algorithmic}
    \Procedure{Make-Set}{$x$}
        \State $x.p = x$
    \EndProcedure
\end{algorithmic}

\subsection{Find-Set}
\begin{algorithmic}
    \Procedure{Find-Set}{$x$}
        \If{$x.p = x$}
            \State \Return $x.p$

            \Else
                \State \Return \Call{Find-Set}{$x.p$}
        \EndIf
    \EndProcedure
\end{algorithmic}


\subsection{Union}
La radice di un albero punta la radice dell'altro

\begin{algorithmic}
    \Procedure{Union}{x,y}
        \State $z = \Call{Find-Set}{$x$}$
        \State $w = \Call{Find-Set}{y}$
        \State $z.p = w$
    \EndProcedure
\end{algorithmic}

Una sequenza di $n-1$ Union può creare una catena lineare di $n$ nodi

\begin{figure}[htbp]
\centering
\begin{tikzpicture}[
    vertex/.style={
        circle,
        draw=black,
        line width=1.3pt,
        minimum size=8.5mm,
        inner sep=0pt,
        font=\large\bfseries
    },
    arrow edge/.style={
        -{Stealth[length=3mm, width=2.2mm]},
        line width=1.3pt,
        draw=black,
        shorten <=2pt,
        shorten >=2pt
    },
    loop edge/.style={
        -{Stealth[length=3mm, width=2.2mm]},
        line width=1.3pt,
        draw=black,
        looseness=6,
        in=100,
        out=25
    },
    op label/.style={
        anchor=west,
        font=\large\bfseries
    }
]

    % Parametri di spaziatura
    % x: a=0, b=2.2, c=4.4, d=6.6
    % y: Riga 1 = 0, Riga 2 = -2.2, Riga 3 = -4.4, Riga 4 = -6.6

    % ==================== RIGA 1 (Stato Iniziale) ====================
    \node[vertex] (a1) at (0, 0) {a};
    \node[vertex] (b1) at (2.2, 0) {b};
    \node[vertex] (c1) at (4.4, 0) {c};
    \node[vertex] (d1) at (6.6, 0) {d};

    \draw[loop edge] (a1) to (a1);
    \draw[loop edge] (b1) to (b1);
    \draw[loop edge] (c1) to (c1);
    \draw[loop edge] (d1) to (d1);

    \node[anchor=west, font=\LARGE\bfseries] at (8.0, 0) {$n=4$};

    % ==================== RIGA 2: UNION(a, b) ====================
    \node[op label] at (-0.6, -1.3) {\textsc{Union}\,(a, b)};

    \node[vertex] (a2) at (0, -2.2) {a};
    \node[vertex] (b2) at (2.2, -2.2) {b};
    \node[vertex] (c2) at (4.4, -2.2) {c};
    \node[vertex] (d2) at (6.6, -2.2) {d};

    \draw[arrow edge] (a2) -- (b2);
    \draw[loop edge] (b2) to (b2);
    \draw[loop edge] (c2) to (c2);
    \draw[loop edge] (d2) to (d2);

    % ==================== RIGA 3: UNION(a, c) ====================
    \node[op label] at (-0.6, -3.5) {\textsc{Union}\,(a, c)};

    \node[vertex] (a3) at (0, -4.4) {a};
    \node[vertex] (b3) at (2.2, -4.4) {b};
    \node[vertex] (c3) at (4.4, -4.4) {c};
    \node[vertex] (d3) at (6.6, -4.4) {d};

    \draw[arrow edge] (a3) -- (b3);
    \draw[arrow edge] (b3) -- (c3);
    \draw[loop edge] (c3) to (c3);
    \draw[loop edge] (d3) to (d3);

    % ==================== RIGA 4: UNION(a, d) ====================
    \node[op label] at (-0.6, -5.7) {\textsc{Union}\,(a, d)};

    \node[vertex] (a4) at (0, -6.6) {a};
    \node[vertex] (b4) at (2.2, -6.6) {b};
    \node[vertex] (c4) at (4.4, -6.6) {c};
    \node[vertex] (d4) at (6.6, -6.6) {d};

    \draw[arrow edge] (a4) -- (b4);
    \draw[arrow edge] (b4) -- (c4);
    \draw[arrow edge] (c4) -- (d4);
    \draw[loop edge] (d4) to (d4);

\end{tikzpicture}
\end{figure}

\newpage
\section{Euristica dell'unione per rango + compressione dei cammini}
\begin{itemize}
    \item \textbf{Altezza di un nodo $\mathbf{x}$}: Numero di archi su un più lungo cammino da $x$ ad una foglia
    \item \textbf{Rango di un nodo $\mathbf{x}$}: Limite superiore dell'altezza di $x$
\end{itemize}

\subsection{Unione per rango}
La radice di rango <punta a quella di rango>.

In tal modo il rango della radice dell'albero così ottenuto non aumenta rispetto al rango, a meno che i ranghi non siano uguali ed in tal caso aumenta di 1.

\subsection{Metodi}

\subsubsection{Make-Set}
\begin{algorithmic}
    \Procedure{Make-Set}{$x$}
        \State $x.p = x$
        \State $x.rank = 0$
    \EndProcedure
\end{algorithmic}

\subsubsection{Union}
\begin{algorithmic}
    \Procedure{Union}{$x,y$}
        \State \Call{Link}{\Call{Find-Set}{$x$},\Call{Find-Set}{$y$}}
    \EndProcedure
\end{algorithmic}

\subsubsection{Link}
\begin{algorithmic}
    \Procedure{Link}{$x,y$}
        \If{$x.rank > y.rank$}
            \State $y.p = x$

            \Else
                \State $x.p = y$

                    \If{$x.rank = y.rank$}
                        \State $y.rank = y.rank + 1$
                    \EndIf
        \EndIf
    \EndProcedure
\end{algorithmic}

\subsubsection{Find-Set con compressione dei cammini}
\begin{algorithmic}
    \Procedure{Find-Set}{$x$}
        \If{$x.p \not = x$}
            \State $x.p = \Call{Find-Set}{$x.p$}$
        \EndIf

        \State \Return $x.p$
    \EndProcedure
\end{algorithmic}





